
\subsubsection{3.1 Overview}

The methodology follows directly from the key insight: benchmark saturation is a structural break, not a monotonic trend. This motivates a three-stage pipeline: (1) remove the global linear trend, (2) detect the structural break in the residuals, (3) validate the break with independent tests.

The full pipeline:
\begin{verbatim}
(i)  Data ingestion: PwC leaderboard → N=115 benchmarks with CoV and paper_count
(ii) OLS detrending: residual_CoV = CoV - (slope × paper_count + intercept)
(iii) PELT detection: structural break at paper_count* from residual CoV series
(iv) Validation: permutation test (p < 0.05) + piecewise F-test (model comparison)
(v)  Regime characterization: F-test (pre vs. global) + Brown-Forsythe (pre vs. post)
\end{verbatim}

\textbf{Rationale for detrending:} Without detrending, PELT would detect the global CoV-paper\_count correlation as a spurious change-point. Detrending isolates the structural break component, making detection robust to the global trend's direction and magnitude (see §5.1 for the OLS trend reversal observation).

\subsubsection{3.2 Data}

The PwC benchmark archive (HuggingFace \texttt{paperswithcode/paperswithcode-data}, August 2026 snapshot) is loaded via Arrow IPC. For each benchmark, CoV = $\sigma$\_b / $\mu$\_b over all reported metric values. Benchmarks with fewer than 38 papers are excluded (min\_papers=38), yielding \textbf{N=115 benchmarks} with paper\_count $\in$ [38, 352].

\subsubsection{3.3 OLS Detrending}

The global CoV-paper\_count relationship confounds change-point detection and is removed via OLS:

residual\_CoV\_b = CoV\_b $-$ ($\beta$̂$_0 + \beta$̂$_1 \cdot$ paper\_count\_b)

This removes whatever linear trend exists (positive or negative) without assuming its direction, ensuring trend-agnostic structural break detection.

\subsubsection{3.4 PELT Change-Point Detection}

PELT [Killick2012] is applied with an L2 cost function (mean-shift detection). The BIC-based penalty $\sigma ^2 \cdot$ log(N) selects the number of breakpoints. Penalties are swept across [1, 50] on a log scale (20 values), confirming a single breakpoint across the relevant range (see Figure 9 in §5).

\textbf{Key parameters:}

\begin{table}[htbp]
\centering
\resizebox{\columnwidth}{!}{%
\begin{tabular}{lll}
\toprule
Parameter & Value & Rationale \\
\midrule
model & L2 & Mean-shift in 1D continuous signal \\
min\_size & 3 & Per VPWBS theoretical recommendation [Wang2019] \\
jump & 1 & Exact search for N=115 \\
BIC penalty & $\sigma ^{2}\cdot log(N) \approx$ 4.32 & Automatic; consistent with PELT theory \\
N\_permutations & 1000 & Standard for permutation test power \\
N\_bootstrap & 1000 & Standard for CI estimation \\
seed & 42 & Fixed for reproducibility \\
\bottomrule
\end{tabular}
}
\end{table}

\subsubsection{3.5 Statistical Validation}

\textbf{Permutation test:} paper\_count labels are shuffled 1000 times; PELT is rerun on each permuted series. The permutation p-value is defined as the fraction of permutations with a detected breakpoint index $\leq$ the observed index (8). This left-tailed formulation tests whether the observed break falls unusually early in the sorted paper\_count distribution under the null of no structural pattern --- a position test specifically targeting early saturation onset. The complementary piecewise F-test (p=0.0021) provides position-agnostic model comparison evidence that does not depend on break position. Gate: permutation p < 0.05.

\textbf{Piecewise F-test (corroboration):} A two-segment linear model (fitted separately on pre/post segments) is compared to a single linear model. A significant F-test confirms the two-segment model explains substantially more variance than the one-segment model.

\subsubsection{3.6 Regime Characterization}

\textbf{H-M1 (exploration regime):} F-test comparing pre-breakpoint variance (n\_pre=8) to global variance. Gate: variance\_ratio > 1.0 AND F-test one-tailed p < 0.10.

\textbf{H-M2 (saturation regime):} Brown-Forsythe test comparing pre-breakpoint variance to post-breakpoint variance. Gate: ratio < 1.0 AND BF p < 0.05.

\textbf{H-M3 (directional concentration):} Four directional metrics evaluated. Gate: $\geq$ 2/4 pass at p < 0.10 (SHOULD\_WORK level).

\subsubsection{3.7 Implementation}

Python 3.10+, \texttt{ruptures} (PELT), \texttt{scipy} (permutation/BF tests), \texttt{statsmodels} (OLS). Data loading via Arrow IPC. Full pipeline automated in \texttt{run\textbackslash \_experiment.py}; hyperparameters centralized in \texttt{config.py}.

